\tongjisetup{
  %******************************
  % 注意：
  %   1. 配置里面不要出现空行
  %   2. 不需要的配置信息可以删除
  %   3. 二级学科我们是计算机应用技术
  %******************************
  % 秘级
  secretlevel={公开},
  secretyear={0},
  %
  %
  % 中文信息
  ctitle={中文标题~English Title, \\ 中文副标题~English SubTitle}, % 可手动加入换行符 \\ 控制换行
  cauthor={姓名},
  studentnumber={1234567},
  cdepartment={电子与信息工程学院},
  csubjectcategory={工学}, % 学科门类
  cmajorfirst={一级学科 \& 专业学位类别}, % 专业学位类别
  cmajorsecond={二级学科 \& 专业领域}, % 专业领域
  cresearchfield={研究方向}, % 研究方向
  csupervisor={导师姓名 教授}, 
  cassosupervisor={副导师姓名 教授}, % 副指导老师/行业导师, 没有把{}中内容留空即可, 或者注释掉
  cjointinstitution={联合培养单位名称}, % 联合培养单位, 没有把{}中内容留空即可, 或者注释掉
  cdate={\zhdigits{2025}年\zhnumber{3}月}, % 手动指定日期, 否则自动使用当前时间
  cfunds={（本论文由国家杰出青年基金 (No.123456789) 支持）}, % 中文版基金信息
  % 英文信息
  etitle={Thesis Title in English}, 
  eauthor={Nome Cognome},
  edepartment={College of Electronic and Information Engineering},
  esubjectcategory={Engineering}, % 学科门类
  emajorfirst={First-level Discipline \& Degree}, % 专业学位类别
  emajorsecond={Second-level Discipline \& Degree's Field}, % 专业领域
  eresearchfield={Research Fieldd}, % 研究方向
  esupervisor={Prof. Supervisor},
  eassosupervisor={Prof. Associate Supervisor}, % 副指导老师/行业导师
  ejointinstitution={Joint Training Institution}, % 联合培养单位
  edate={March, 2025}, % 手动指定日期, 否则自动使用当前时间
  efunds={(Supported by the Natural Science Foundation of China for Distinguished Young Scholars, Grant No.123456789)} % 英文版基金信息
}

% 后处理
\makeatletter
\define@key{tongji}{cheadingtitlenew}{\cheadingtitlenew{#1}}
\gdef\cheadingtitlenew#1{\gdef\tongji@cheadingtitlenew{#1}}
\tongjideletelinebreakcmd{tongji@ctitle}{tongji@posttitle}
\cheadingtitlenew{\tongji@posttitle}
\makeatother

% 定义中英文摘要和关键字
% 背景+(1)(2)(3)三个工作内容，123分别讲述针对什么问题，做了什么工作，解决了什么问题。
\begin{cabstract}
  \LaTeX~是常用的排版软件之一，能够输出高质量的学术出版物。
  \tongjithesis{}~是基于~\LaTeX~的模板，历经多次改进。本次改进的主要贡献如下：
  
  \begin{enumerate}[1.]
    \item 将平台迁移到 \emph{vscode}；
    \item 适配最新的同济大学硕士学位论文要求；
    \item 添加了包括 \emph{tikz}、\emph{algorithm2e} 等宏包的支持。
  \end{enumerate}

  感谢使用！
\end{cabstract}

\ckeywords{LaTeX, 毕业论文, 模板, 同济大学}

% 后处理
\makeatletter
\define@key{tongji}{cheadingtitlenew}{\cheadingtitlenew{#1}}
\gdef\cheadingtitlenew#1{\gdef\tongji@cheadingtitlenew{#1}}
\tongjideletelinebreakcmd{tongji@ctitle}{tongji@posttitle}
\cheadingtitlenew{\tongji@posttitle}
% === 新增：将英文摘要的页眉强制改为 Abstract ===
\patchcmd{\tongji@makeabstract}
  {\xparameter@heading@eother{\tongji@eabstractname}}
  {\xparameter@heading@eother{Abstract}}
  {}{}
% ===============================================
\makeatother

% ===============================================
% === 新增：生成单独书脊页的命令 ===
\newcommand{\makespine}{
  \clearpage
  \thispagestyle{empty}
  % 页面总高 29.7cm，减去上下各 5cm，中间文本区高度设为 19.7cm。
  % 利用 \vspace*{\fill} 上下弹性挤压，完美实现物理边距各 5cm。
  \vspace*{\fill}
  \begin{center}
    % 宽度设为 1.2em，强制每个汉字自动换行，实现垂直排版
    \parbox[c][19.7cm][s]{1.2em}{
      \centering
      \fangsong\zihao{4}\bfseries         % 仿宋、四号、加粗
      \fontsize{14bp}{16bp}\selectfont    % 四号字对应14bp，行距严格指定为16bp(磅)
      \setlength{\parskip}{0pt}           % 段前段后 0 磅
      
      \tongji@posttitle                   % 使用无换行符版本的标题
      退化条件下的鲁棒语音转换算法研究
      \par\vfill
      \tongji@cauthor  
      宋昕宁                              % 姓名
      \par\vfill
      同济大学
    }
  \end{center}
  \vspace*{\fill}
  \clearpage
}
% ===============================================
\makeatother

\begingroup
\ctexset{
  chapter={
    % format 设置：单倍行距、居中、Arial字体(sffamily)、三号字(zihao{3})、加粗(bfseries)
    format={\linespread{1}\selectfont\centering\sffamily\zihao{3}\bfseries},
    % beforeskip={24bp}, % 段前 24 磅 (LaTeX中的 bp 对应 Word 的磅)
    afterskip={18bp}   % 段后 18 磅
  }
}

\begin{eabstract}
  \LaTeX~is the most popular typesetting system for academic writing.
  \tongjithesis{}~is a~\LaTeX~template for Tongji University. The main contributions of this modification are as follows:

  \begin{enumerate}[1.]
    \item Migrate the platform to \emph{vscode};
    \item Adapt to the latest requirements of the master's thesis of Tongji University;
    \item Add support for packages such as \emph{tikz} and \emph{algorithm2e}.
  \end{enumerate}

  Thanks for using!
\end{eabstract}

\ekeywords{LaTeX, thesis, template, tongji-university}
